\documentclass[a4paper,11pt]{article}
\usepackage[a4paper,margin=25mm]{geometry}
\usepackage{amsmath,amssymb,amsfonts}
\usepackage{hyperref}
\usepackage{booktabs}

\title{Penurunan Persamaan Simulator Magnetoimpedansi Teoretis\\dengan Sensing Nanopartikel\\{\normalsize Penjelasan verbose langkah-demi-langkah (versi revisi)}}
\author{Disusun dari \texttt{simulator-magnetoimpedansi-teoretis-revisi.html}}
\date{\today}

\begin{document}
\maketitle

\begin{abstract}
Dokumen ini menulis ulang seluruh teori yang dipakai simulator dengan bahasa
yang dibuat semudah mungkin untuk konfirmasi ulang.
Penomoran dibuat otomatis oleh LaTeX (lingkungan \texttt{equation} plus
\texttt{subequations} untuk kelompok berhuruf), sehingga urutannya
mengikuti panel teori pada simulator (23 nomor utama),
ditambah semua ekstensi empiris pada kode:
koreksi tepi, dinding domain, sebaran koersif, sebaran anisotropi,
sebaran diameter, konduktivitas partikel, cakupan lapisan,
lapisan dua sisi, skin kabel, kalibrasi kumparan, derau,
pita model satu sigma, dan batas deteksi.
Semua hubungan matematis ditulis pada baris khusus
(display equation), tidak ada persamaan yang disembunyikan
di dalam paragraf.
Konvensi yang dipakai adalah konvensi fasor waktu maju,
satuan SI, dan medan bias dinyatakan sebagai rapat fluks
dalam tesla (mili-tesla pada tampilan).
\end{abstract}

\tableofcontents
\newpage

%=====================================================================
\section{Cara membaca dokumen ini}
%=====================================================================
Paragraf teks hanya berisi cerita dan alasan fisis.
Setiap kali ada hubungan antar besaran, hubungan itu
ditulis pada baris tersendiri di bawah paragrafnya.
Jadi pembaca dapat memeriksa satu persamaan dalam satu waktu,
mencocokkannya dengan satu fungsi pada kode,
lalu lanjut ke langkah berikutnya.

Pemetaan kasar fungsi pada kode:
fungsi permeabilitas adalah \texttt{chiLayer},
fungsi impedansi trilayer adalah \texttt{zOne},
rata-rata histeresis adalah \texttt{zSensor},
rata-rata sebaran anisotropi adalah \texttt{zSensorAvg},
respons partikel adalah \texttt{particleResp} dan \texttt{particleAvg},
geometri lapisan adalah \texttt{layerGeom},
impedansi total adalah \texttt{totalZ},
sapuan dan derau adalah \texttt{setpoints}, \texttt{measure},
\texttt{modelCurve}, \texttt{peakAt}, \texttt{sigmaMIR},
\texttt{modelBand}, dan \texttt{modelMIR}.

Rujukan silang memakai mekanisme standar LaTeX.
Setiap persamaan bernomor memiliki \texttt{\textbackslash label\{eq:...\}},
misalnya Helmholtz \texttt{eq:1}, difusi \texttt{eq:3a},
MIR \texttt{eq:2}, dan batas deteksi \texttt{eq:23b}.
Untuk merujuk, pakai \texttt{\textbackslash eqref\{eq:1\}} yang menghasilkan
\eqref{eq:1}, atau \texttt{\textbackslash ref\{eq:3a\}} sesuai kebutuhan.
Kelompok berhuruf memakai lingkungan \texttt{subequations} dengan label
grup (\texttt{eq:3}, \texttt{eq:4}, dan seterusnya), sehingga penomoran
huruf dibuat otomatis dan tidak perlu \texttt{\textbackslash tag}
manual.

%=====================================================================
\section{Catatan revisi dan status tiap bagian}
%=====================================================================
Versi ini memperbaiki kesalahan yang ditemukan pada tinjauan sebelumnya,
baik pada dokumen maupun pada kode simulator.
Kolom kedua menunjukkan di mana perbaikan dilakukan.
\begin{center}
\small
\begin{tabular}{p{0.30\textwidth}p{0.12\textwidth}p{0.48\textwidth}}
\toprule
Butir & Lokasi & Perubahan\\
\midrule
Rugi eddy bola logam & kode dan dokumen & Faktor dua diperbaiki;
bentuk akhir pada sub-bab rugi eddy.\\
Dimensi magnetisasi para/dia & dokumen & Definisi dibetulkan
(lihat sub-bab para/dia); kode sudah benar.\\
Demagnetisasi film sangat tipis & dokumen & Dibetulkan
(lihat sub-bab demagnetisasi); kode sudah benar.\\
Nilai faktor amplitudo acak & dokumen & Dibetulkan
(lihat sub-bab faktor amplitudo).\\
Simpul Gauss--Hermite & dokumen & Simpul fisikawan \emph{dikali} $\sqrt{2}$.\\
Tanda $H_{y}$ & dokumen & Tanda minus Faraday dipertahankan.\\
Penomoran persamaan & dokumen & Persamaan bagian notasi tidak lagi bernomor.\\
Notasi bertabrakan & dokumen & $\mathcal{A},\mathcal{B}$, $\Delta_{c}$,
$m_{p}$, $\varpi$ dibedakan dari simbol lain; $\alpha_{p}$ didefinisikan.\\
Parser CSV & kode & Fallback yang salah ke kolom $I_{\mathrm{aktual}}$ dihapus.\\
Suseptibilitas partikel & kode dan dokumen & Terblokir memakai
$\chi_{\perp,b}$ transversal, konsisten dengan superparamagnetik.\\
Sebaran diameter & kode dan dokumen & Rata-rata berbobot volume ($d^{3}$).\\
Koreksi tepi & kode dan dokumen & Bentuk aditif $\pm N_{w}\mu_{0}M_{s}$,
hanya tebal lapisan magnetik.\\
Dinding domain & kode dan dokumen & Bentuk Debye kausal (real dan imajiner).\\
Acuan MIR & kode dan dokumen & Diambil pada medan \emph{aktual} di ujung sapuan.\\
Sensitivitas LOD & kode dan dokumen & Langkah $dm$ adaptif pada rezim linear.\\
Klaim ``sudah diuji'' & kode dan dokumen & Ada fungsi uji mandiri; hasil
pada Lampiran~\ref{app:uji}.\\
\bottomrule
\end{tabular}
\end{center}
Beberapa kelemahan fisis bersifat struktural dan \emph{tidak} diubah,
tetapi sekarang ditandai terang-terangan (Lampiran~\ref{app:empiris}):
medan sebar seragam, sebaran dipolar Gaussian, remanen model tanh,
derau acuan common-mode, rata-rata naik-turun pada data campuran arah,
dan seluruh parameter heuristik.

%=====================================================================
\section{Alur perhitungan ujung-ke-ujung dan peta hubungan rumus}
%=====================================================================
Bagian ini adalah benang merah dokumen.
Tujuannya agar pembaca selalu tahu posisi setiap rumus:
datang dari mana, dipakai untuk apa, dan mengalir ke mana.
Bagian-bagian sesudahnya adalah penurunan rinci tiap kotak pada alur ini.

\subsection{Gambaran satu kalimat}
Pengguna meminta medan.
Kumparan mengubah permintaan menjadi medan aktual.
Massa partikel mengubah medan itu menjadi tiga efek.
Sensor mengubah medan efektif menjadi impedansi.
Semua kontribusi dijumlahkan menjadi impedansi total,
lalu dinormalisasi terhadap ujung sapuan menjadi persen MIR,
lalu dibandingkan dengan data berderau melalui pita satu sigma.

\subsection{Diagram pipa sepuluh tahap}
Tahap nol menyiapkan parameter.
Semua geseran dan angka bahan dibaca dari panel dan dikemas menjadi
struktur parameter sensor dan struktur bahan partikel.
Tidak ada fisika pada tahap ini, hanya penyiapan.

Tahap satu mengubah keinginan menjadi kenyataan medan.
Masukan tahap ini adalah medan yang diinginkan.
Keluaran tahap ini adalah tiga medan: perintah, aktual, dan terbaca.
Arus perintah diperoleh sebagai berikut.
\begin{equation*}
I_{\mathrm{cmd}} = H_{\mathrm{set}}/k_{\mathrm{cmd}}
\end{equation*}
Medan aktual diperoleh sebagai berikut.
\begin{equation*}
H_{\mathrm{act}} = H_{\mathrm{set}}\,(k_{\mathrm{true}}/k_{\mathrm{cmd}})
+ B_{\mathrm{off}} + \delta H
\end{equation*}
Medan terbaca diperoleh sebagai berikut.
\begin{equation*}
H_{\mathrm{read}} = H_{\mathrm{act}} + \delta H_{r}
\end{equation*}
Hanya medan aktual yang masuk ke fisika.
Dua lainnya hanya untuk sumbu plot, arus, dan berkas CSV.
Rujukan rinci ada pada \eqref{eq:1}.

Tahap dua mengubah massa menjadi geometri lapisan.
Masukan tahap ini adalah massa, massa jenis, fraksi volume, luas sebaran,
dan lebar strip.
Keluaran tahap ini adalah tebal, tebal efektif, dan faktor demagnetisasi.
Tebal mengikuti definisi massa dibagi volume sebagai berikut
(lihat \eqref{eq:17a}).
\begin{equation*}
t_{p} \Leftarrow m/(\rho\,\varphi\,A_{c})
\end{equation*}
Tebal efektif jenuh pada skala lebar sebagai berikut
(lihat \eqref{eq:17b}, \eqref{eq:17c}).
\begin{equation*}
t_{\mathrm{eff}} \Leftarrow \lambda\,(1-e^{-t_{p}/\lambda})
\end{equation*}
Faktor demagnetisasi mengikuti geometri pipih sebagai berikut.
\begin{equation*}
N_{\mathrm{in}} \Leftarrow (1-N_{z})/2
\end{equation*}
Panah kiri dibaca ``dihitung dari''.
Fungsi kode tahap ini adalah \texttt{layerGeom} dan \texttt{Nin}.

Tahap tiga mengubah medan menjadi respons partikel.
Masukan tahap ini adalah medan dalam tesla, arah sapuan, bahan, frekuensi,
dan suhu.
Keluaran tahap ini adalah magnetisasi partikel, faktor amplitudo momen,
dan suseptibilitas kompleks.
Untuk superparamagnetik, rantainya adalah sebagai berikut
(lihat \eqref{eq:12a}--\eqref{eq:14b}).
\begin{equation*}
B \Rightarrow \xi \Rightarrow L \Rightarrow M \Rightarrow \chi_{\perp}
\Rightarrow \chi(\omega)
\end{equation*}
Untuk terblokir, rantainya adalah sebagai berikut
(lihat \eqref{eq:15a}--\eqref{eq:16c}).
\begin{equation*}
B \Rightarrow a \Rightarrow M,\qquad
B \Rightarrow \chi_{\perp,b} \Rightarrow \chi(\omega)
\end{equation*}
Panah kanan dibaca ``menghasilkan''.
Rugi eddy logam ditambahkan ke bagian imajiner bila konduktivitas positif.
Sebaran diameter merata-ratakan ketiga keluaran di atas
dengan kuadratur yang menjaga rata-rata.
Fungsi kode tahap ini adalah \texttt{particleResp} dan \texttt{particleAvg}.

Tahap empat mengubah medan efektif menjadi impedansi sensor telanjang.
Masukan tahap ini adalah medan efektif, arah sapuan, parameter sensor,
dan frekuensi.
Keluaran tahap ini adalah impedansi sensor sebelum lapisan partikel.
Rantainya adalah sebagai berikut
(lihat \eqref{eq:9}--\eqref{eq:11} lalu \eqref{eq:3a}--\eqref{eq:6b}).
\begin{equation*}
H \Rightarrow (A,B) \Rightarrow \chi \Rightarrow \mu
\Rightarrow (k_{c},k_{f}) \Rightarrow (\mathcal{A},\mathcal{B})_{\mathrm{sambung}}
\Rightarrow (E_{s},H_{s}) \Rightarrow Z
\end{equation*}
Dua averaging membungkus rantai ini.
Yang pertama meratakan dua cabang magnetisasi untuk histeresis koersif.
Yang kedua meratakan sebaran anisotropi film.
Fungsi kode tahap ini adalah \texttt{chiLayer}, \texttt{zOne},
\texttt{zSensor}, dan \texttt{zSensorAvg}.

Tahap lima menghitung tiga efek partikel dari keluaran tahap dua dan tiga.
Efek dinamis memakai medium efektif sebagai berikut
(lihat \eqref{eq:18a}--\eqref{eq:19}).
\begin{equation*}
(\varphi,\chi) \Rightarrow \mu_{L} \Rightarrow \Delta Z_{p}
\end{equation*}
Efek rata-rata memakai demagnetisasi sebagai berikut
(lihat \eqref{eq:20}).
\begin{equation*}
(N_{\mathrm{in}},\varphi,M_{p}) \Rightarrow B_{s}
\end{equation*}
Efek sebaran memakai lebar dipolar sebagai berikut
(lihat \eqref{eq:21a}--\eqref{eq:21b}).
\begin{equation*}
(C_{d},\varphi,M_{s},s) \Rightarrow \sigma_{s} \Rightarrow \bar Z
\end{equation*}
Cakupan pulau mengecilkan ketiganya, lapisan dua sisi menggandakan
suku dinamis saja.

Tahap enam menjumlahkan semuanya menjadi impedansi total
(lihat \eqref{eq:22}).
\begin{equation*}
Z_{\mathrm{tot}} = \bar Z + \Delta Z_{p} + R_{c} + j\,\omega\,L_{p}
\end{equation*}
Resistansi kabel memakai bentuk akar frekuensi sebagai berikut.
\begin{equation*}
R_{c}(f) = R_{c0} + R_{c,\mathrm{skin}}\,\sqrt{f_{\mathrm{MHz}}}
\end{equation*}
Modulus diambil sebagai berikut.
\begin{equation*}
|Z| = \sqrt{(\mathrm{Re}\,Z)^{2}+(\mathrm{Im}\,Z)^{2}}
\end{equation*}
Fungsi kode tahap ini adalah \texttt{totalZ}.

Tahap tujuh menormalisasi menjadi persen MIR
(lihat \eqref{eq:2}).
Acuan tiap massa diambil di ujung naik sebagai berikut.
\begin{equation*}
Z_{\mathrm{ref}} = |Z(H_{\max}^{\mathrm{akt}},\mathrm{naik})|
\end{equation*}
Kurva model diperoleh sebagai berikut.
\begin{equation*}
y(H) = 100\,(|Z(H)|-Z_{\mathrm{ref}})/Z_{\mathrm{ref}}
\end{equation*}

Tahap delapan menambahkan derau pengukuran.
Acuan ikut berderau, tiap titik ikut berderau medan dan modulus,
bacaan probe ikut berderau sumbu.
Bentuknya ada pada bab sapuan, intinya adalah sebagai berikut:
derau medan masuk fisika, derau baca hanya menggeser titik.
Fungsi kode tahap ini adalah \texttt{measure}.

Tahap sembilan menilai hasil.
Puncak, sensitivitas awal, dan batas deteksi mengikuti
\eqref{eq:23a}--\eqref{eq:23b}.
Pita model, residu, dan biaya fit memakai satu definisi sigma yang sama,
yaitu gabungan turunan kurva kali derau medan plus derau impedansi.
Fungsi kode tahap ini adalah \texttt{peakAt}, \texttt{sigmaMIR},
\texttt{slopeAt}, \texttt{modelBand}, \texttt{modelMIR},
dan \texttt{nelderMead}.

\subsection{Peta hubungan: tiga aliran bertemu di satu titik}
Ada tiga aliran yang akhirnya bertemu.
Aliran sensor membawa permeabilitas ke impedansi.
Isinya adalah \eqref{eq:8}--\eqref{eq:11} yang mengalir ke
\eqref{eq:3a}--\eqref{eq:6b}.
Aliran partikel membawa magnetisasi ke suseptibilitas.
Isinya adalah \eqref{eq:12a}--\eqref{eq:16c} plus eddy dan lognormal.
Aliran geometri membawa massa ke tebal dan demagnetisasi.
Isinya adalah \eqref{eq:17a}--\eqref{eq:17c} plus faktor spheroid.
Ketiga aliran bertemu pada impedansi total \eqref{eq:22}:
aliran sensor memberi $\bar Z$,
aliran partikel plus geometri memberi $\Delta Z_{p}$, $B_{s}$, $\sigma_{s}$,
lalu $B_{s}$ dan $\sigma_{s}$ dibalikkan lagi ke aliran sensor
sebagai medan efektif.
Inilah alasan mengapa urutan hitung pada kode adalah:
geometri dulu, partikel dulu, baru sensor dirata-rata,
baru dijumlahkan.

\subsection{Jejak satu titik dari awal sampai persen}
Bagian ini menelusuri satu titik sapuan dari permintaan sampai persen,
agar hubungan antar rumus terlihat sebagai satu cerita.
Angka tidak ditulis karena yang penting adalah urutannya.

Pertama, permintaan diubah menjadi arus perintah.
\begin{equation*}
\text{permintaan} \Rightarrow \text{arus perintah}
\end{equation*}
Kedua, arus perintah diubah menjadi medan aktual lewat kalibrasi.
\begin{equation*}
\text{arus perintah} \Rightarrow \text{medan aktual}
\end{equation*}
Ketiga, medan aktual diubah menjadi bacaan probe.
\begin{equation*}
\text{medan aktual} \Rightarrow \text{medan terbaca}
\end{equation*}
Keempat, massa diubah menjadi geometri lapisan.
\begin{equation*}
\text{massa} \Rightarrow \text{tebal} \Rightarrow \text{tebal efektif}
\Rightarrow \text{demagnetisasi}
\end{equation*}
Kelima, medan aktual diubah menjadi respons partikel.
\begin{equation*}
\text{medan aktual} \Rightarrow \text{magnetisasi}
\Rightarrow \text{suseptibilitas}
\end{equation*}
Keenam, respons dan geometri diubah menjadi tiga efek.
\begin{equation*}
\text{respons + geometri} \Rightarrow
(\Delta Z_{p},\,B_{s},\,\sigma_{s})
\end{equation*}
Ketujuh, medan aktual dikoreksi menjadi medan efektif sensor.
\begin{equation*}
\text{medan aktual + sebaran rata-rata} \Rightarrow \text{medan efektif}
\end{equation*}
Kedelapan, medan efektif plus sebaran dipolar diubah menjadi
impedansi sensor rata-rata.
\begin{equation*}
\text{medan efektif + sebaran} \Rightarrow \bar Z
\end{equation*}
Kesembilan, semua suku dijumlahkan dan diambil modulusnya.
\begin{equation*}
\text{jumlah suku} \Rightarrow |Z|
\end{equation*}
Kesepuluh, modulus dinormalisasi terhadap ujung menjadi persen.
\begin{equation*}
\text{modulus + acuan ujung} \Rightarrow \text{persen MIR}
\end{equation*}
Kesebelas, persen model dibandingkan dengan persen terukur berderau
melalui pita, residu, dan khi-kuadrat.
\begin{equation*}
\text{model + data} \Rightarrow \text{pita, residu, kecocokan}
\end{equation*}

\subsection{Tabel ringkas fungsi lawan persamaan}
Tabel ini dipakai sebagai indeks cepat saat membaca kode berdampingan
dengan dokumen.
\begin{center}
\begin{tabular}{lll}
\toprule
Fungsi kode & Persamaan & Peran dalam alur\\
\midrule
\texttt{chiLayer} & \eqref{eq:8}--\eqref{eq:11} & Medan menjadi permeabilitas\\
\texttt{zOne} & \eqref{eq:3a}--\eqref{eq:6b} & Permeabilitas menjadi impedansi\\
\texttt{zSensor} & campuran dua cabang & Histeresis koersif\\
\texttt{zSensorAvg} & rata-rata tujuh titik & Sebaran anisotropi\\
\texttt{particleResp} & \eqref{eq:12a}--\eqref{eq:16c} & Medan menjadi respons partikel\\
\texttt{particleAvg} & rata-rata tujuh titik & Sebaran diameter\\
\texttt{layerGeom}, \texttt{Nin} & \eqref{eq:17a}--\eqref{eq:17c} & Massa menjadi geometri\\
\texttt{totalZ} & \eqref{eq:18a}--\eqref{eq:22} & Pertemuan tiga aliran\\
\texttt{measure} & rantai kalibrasi + derau & Model menjadi data terukur\\
\texttt{sigmaMIR}, \texttt{modelBand} & \eqref{eq:23a} & Data menjadi pita\\
\texttt{peakAt}, LOD & \eqref{eq:23a}--\eqref{eq:23b} & Kurva menjadi batas deteksi\\
\bottomrule
\end{tabular}
\end{center}

%=====================================================================
\section{Notasi, satuan, dan konstanta}
%=====================================================================
Bagian ini mendefinisikan semua simbol agar tidak ada ambiguitas
antara milimeter, mikrometer, mega-hertz, dan tesla.

Permeabilitas vakum didefinisikan sebagai berikut.
\begin{equation*}
\mu_{0} = 4\pi\times 10^{-7}\ \mathrm{H/m}
\end{equation*}
Konstanta Boltzmann didefinisikan sebagai berikut.
\begin{equation*}
k_{B} = 1.380649\times 10^{-23}\ \mathrm{J/K}
\end{equation*}
Rasio giromagnetik didefinisikan sebagai berikut.
\begin{equation*}
\gamma = 1.76085963\times 10^{11}\ \mathrm{rad/(s\,T)}
\end{equation*}
Frekuensi sudut didefinisikan dari frekuensi biasa sebagai berikut.
\begin{equation*}
\omega = 2\pi f
\end{equation*}
Frekuensi yang dinormalisasi terhadap giromagnetik adalah sebagai berikut.
\begin{equation*}
\omega_{\gamma} = \omega / \gamma
\end{equation*}

Geometri strip dijelaskan sebagai berikut.
Panjang strip sepanjang arus adalah $l$.
Lebar strip adalah $w$.
Tebal tiap lapis magnetik adalah $t_{f}$.
Tebal total inti tembaga adalah $t_{\mathrm{cu}}$.
Setengah tebal inti tembaga adalah $d_{c}$.
Hubungan tebal inti adalah sebagai berikut.
\begin{equation*}
t_{\mathrm{cu}} = 2d_{c}
\end{equation*}

Besaran magnetik dijelaskan sebagai berikut.
Simbol $M_{s}$ adalah magnetisasi jenuh.
Simbol $H_{k}$ adalah medan anisotropi.
Simbol $H_{c}$ adalah medan koersif sensor.
Simbol $\alpha$ adalah redaman Gilbert.
Simbol $\Delta H_{0}$ adalah pelebaran garis inhomogen.
Simbol $\sigma_{f}$ adalah konduktivitas lapis magnetik.
Simbol $\sigma_{c}$ adalah konduktivitas tembaga (huruf kecil, sama dengan persamaan di bawah).
Pada tampilan, konduktivitas dalam MS/m sehingga berlaku konversi berikut.
\begin{equation*}
\sigma[\mathrm{S/m}] = \sigma[\mathrm{MS/m}]\times 10^{6}
\end{equation*}
Tebal film dalam mikrometer sehingga berlaku konversi berikut.
\begin{equation*}
t[\mathrm{m}] = t[\mu\mathrm{m}]\times 10^{-6}
\end{equation*}
Panjang dan lebar dalam milimeter sehingga berlaku konversi berikut.
\begin{equation*}
l[\mathrm{m}] = l[\mathrm{mm}]\times 10^{-3}
\end{equation*}
\begin{equation*}
w[\mathrm{m}] = w[\mathrm{mm}]\times 10^{-3}
\end{equation*}
Medan pada kode dalam mili-tesla sehingga berlaku konversi berikut.
\begin{equation*}
H[\mathrm{T}] = H[\mathrm{mT}]\times 10^{-3}
\end{equation*}
Catatan: besaran yang disebut ``medan'' $H$ pada dokumen dan kode sebenarnya
rapat fluks $\mu_{0}H$ dalam tesla. Karena itu $H_{k}$, $H_{c}$, dan
$\mu_{0}M_{s}$ semuanya memiliki satuan yang sama dan dapat dijumlahkan
langsung.

%=====================================================================
\section{Medan kumparan Helmholtz dan definisi MIR}
%=====================================================================
\subsection{Medan di pusat kumparan}
Dua kumparan lingkaran identik dengan jari-jari $R$,
diletakkan sejajar pada jarak $R$,
masing-masing dengan $N$ lilitan dan dialiri arus $I$ yang sama.
Dari hukum Biot--Savart, medan aksial satu loop pada jarak $z$ adalah
sebagai berikut.
\begin{equation*}
B_{\mathrm{loop}}(z) = \frac{\mu_{0}NI R^{2}}{2\left(R^{2}+z^{2}\right)^{3/2}}
\end{equation*}
Pusat Helmholtz berada di tengah, yaitu pada jarak setengah jari-jari
dari tiap loop, yang ditulis sebagai berikut.
\begin{equation*}
z = R/2
\end{equation*}
Menjumlahkan dua loop memberi faktor geometri sebagai berikut.
\begin{equation*}
B = 2\times \frac{\mu_{0}NI R^{2}}{2\left(R^{2}+R^{2}/4\right)^{3/2}}
\end{equation*}
Penyederhanaan faktor geometri memberi bentuk yang dipakai simulator.
\begin{equation}
B = \left(\frac{4}{5}\right)^{3/2}\frac{\mu_{0}NI}{R}
\label{eq:1}
\end{equation}
Nilai numerik faktor geometri adalah sebagai berikut.
\begin{equation*}
\left(\frac{4}{5}\right)^{3/2} \approx 0.7155
\end{equation*}
Sehingga bentuk praktisnya adalah sebagai berikut.
\begin{equation*}
B \approx 0.7155\,\frac{\mu_{0}NI}{R}
\end{equation*}
Kode menghitung konstanta teori dalam mT/A sebagai berikut.
\begin{equation*}
k_{\mathrm{teori}} = 0.7155\,\mu_{0}\,N/(R/100)\times 1000
\end{equation*}
Alasan dibagi 100 adalah karena $R$ pada tampilan dalam cm,
sedangkan rumus SI meminta meter.

\subsection{Rantai kalibrasi: keinginan, kenyataan, bacaan}
Simulator membedakan tiga medan agar kalibrasi dapat diuji.
Pertama, medan yang diinginkan oleh pengguna.
\begin{equation*}
H_{\mathrm{set}} = h
\end{equation*}
Kedua, arus perintah yang dikirim ke catu daya.
\begin{equation*}
I_{\mathrm{cmd}} = H_{\mathrm{set}}/k_{\mathrm{cmd}}
\end{equation*}
Ketiga, medan aktual yang benar-benar dirasakan sensor.
\begin{equation*}
H_{\mathrm{act}} = H_{\mathrm{set}}\,(k_{\mathrm{true}}/k_{\mathrm{cmd}})
+ B_{\mathrm{off}} + \delta H
\end{equation*}
Di sini $k_{\mathrm{cmd}}$ adalah konstanta yang diyakini alat,
$k_{\mathrm{true}}$ adalah konstanta sebenarnya,
$B_{\mathrm{off}}$ adalah offset Bumi plus remanen,
dan $\delta H$ adalah fluktuasi acak medan.
Keempat, medan terbaca oleh Hall probe.
\begin{equation*}
H_{\mathrm{read}} = H_{\mathrm{act}} + \delta H_{r}
\end{equation*}
Derau $\delta H_{r}$ hanya menggeser titik pada sumbu X,
tidak masuk ke fisika impedansi.
Fisika $Z$ selalu memakai $H_{\mathrm{act}}$.
Sumbu plot memakai $H_{\mathrm{set}}$ atau $H_{\mathrm{read}}$
sesuai saklar tampilan.
Arus aktual yang konsisten dengan medan aktual adalah sebagai berikut.
\begin{equation*}
I_{\mathrm{act}} = (H_{\mathrm{act}}-B_{\mathrm{off}})/k_{\mathrm{true}}
\end{equation*}
Bila konstanta nol, kode menjaga agar arus menjadi tak-hingga yang jelas
(Not-a-Number) daripada diam-diam salah.

\subsection{Definisi MIR}
Rasio magnetoimpedansi didefinisikan terhadap modulus di ujung sapuan.
\begin{equation}
\mathrm{MIR}(H) = 100\,\frac{|Z(H)|-|Z(H_{\max})|}{|Z(H_{\max})|}
\label{eq:2}
\end{equation}
Satuan MIR adalah persen.
Yang dimaksud $H_{\max}$ pada definisi ini adalah medan \emph{aktual}
di ujung sapuan, bukan medan yang diminta.
\begin{equation*}
H_{\max}^{\mathrm{akt}} = H_{\max}\,k_{\mathrm{true}}/k_{\mathrm{cmd}}
+ B_{\mathrm{off}}
\end{equation*}
Bila kalibrasi default (tanpa selisih konstanta dan tanpa offset),
medan aktual ujung identik dengan medan nominal.
Alasan memakai ujung sapuan adalah agar definisi tidak memerlukan
ekstrapolasi ke tak-hingga.
Konsekuensinya, bila kurva belum jenuh di $H_{\max}$,
penyebut bias dan puncak ikut bias.
Kode memberi peringatan bila dua titik terakhir masih berubah lebih dari
ambang setengah poin atau dua persen relatif.

%=====================================================================
\section{Impedansi trilayer dari persamaan Maxwell}
%=====================================================================
\subsection{Dari Maxwell ke persamaan difusi}
Arus mengalir sepanjang $x$, semua besaran hanya berubah sepanjang $z$
karena lebar jauh lebih besar daripada tebal.
Hukum Faraday dan Ampere tanpa arus perpindahan adalah sebagai berikut.
\begin{equation*}
\nabla\times \mathbf{E} = -j\omega \mathbf{B}
\end{equation*}
\begin{equation*}
\nabla\times \mathbf{H} = \sigma \mathbf{E}
\end{equation*}
Hubungan konstitutif magnetik adalah sebagai berikut.
\begin{equation*}
\mathbf{B} = \mu_{0}\,\mu\,\mathbf{H}
\end{equation*}
Eliminasi $\mathbf{H}$ memberi persamaan difusi untuk komponen $E_{x}$.
\begin{subequations}\label{eq:3}
\begin{equation}
\frac{d^{2}E_{x}}{dz^{2}} = k^{2}E_{x}
\label{eq:3a}
\end{equation}
Bilangan gelombang kompleks didefinisikan sebagai berikut.
\begin{equation}
k = \sqrt{j\omega\mu_{0}\mu\sigma}
\label{eq:3b}
\end{equation}
Bentuk skin-depth diperoleh dengan menulis akar dari satuan imajiner.
\begin{equation}
k = \frac{1+j}{\delta}
\label{eq:3c}
\end{equation}
Skin-depth didefinisikan sebagai berikut.
\begin{equation}
\delta = \sqrt{\frac{2}{\omega\mu_{0}\mu\sigma}}
\label{eq:3d}
\end{equation}
\end{subequations}
Makin besar frekuensi, konduktivitas, atau permeabilitas,
makin kecil skin-depth dan arus makin menumpuk di permukaan.

\subsection{Solusi simetris pada inti dan lapis magnetik}
Struktur simetris terhadap tengah.
Pada inti tembaga, solusi genap yang memenuhi simetri adalah kosinus
hiperbolik sebagai berikut.
\begin{subequations}\label{eq:4}
\begin{equation}
E_{c}(z) = \cosh(k_{c}z)
\label{eq:4a}
\end{equation}
Pada lapis magnetik, dipakai koordinat lokal dari antarmuka sebagai berikut.
\begin{equation*}
u = z-d_{c}
\end{equation*}
Solusi umum kombinasi kosinus dan sinus hiperbolik adalah sebagai berikut.
\begin{equation}
E_{f}(u) = \mathcal{A}\cosh(k_{f}u)+\mathcal{B}\sinh(k_{f}u)
\label{eq:4b}
\end{equation}
\end{subequations}
Bilangan gelombang tiap bahan adalah sebagai berikut.
\begin{equation*}
k_{c} = \sqrt{j\omega\mu_{0}\sigma_{c}}
\end{equation*}
\begin{equation*}
k_{f} = \sqrt{j\omega\mu_{0}\sigma_{f}\,\mu}
\end{equation*}
Permeabilitas kompleks adalah sebagai berikut.
\begin{equation*}
\mu = 1+\chi
\end{equation*}

\subsection{Koefisien sambung dari syarat batas}
Medan magnetik transversal diperoleh dari hukum Faraday, termasuk
tanda minusnya.
\begin{equation*}
H_{y} = -\frac{1}{j\omega\mu_{0}\mu}\frac{dE_{x}}{dz}
\end{equation*}
Tanda minus itu fisis: untuk arus pada arah $+x$, medan $H_{y}$ di muka atas
($z>0$) bernilai negatif (kaidah tangan kanan). Karena itu besaran yang
dipakai pada hukum Ampere adalah komponen yang positif untuk arus positif.
\begin{equation*}
H_{s} \equiv -H_{y}\big|_{\mathrm{muka\ atas}}
= \frac{1}{j\omega\mu_{0}\mu}\frac{dE_{x}}{dz}\bigg|_{\mathrm{muka\ atas}}
\end{equation*}
Syarat kontinuitas $H_{y}$ di antarmuka tidak terpengaruh tanda ini karena
kedua sisi memuat faktor $-1$ yang sama.
Kontinuitas komponen listrik di antarmuka menuntut nilai yang sama.
Posisi antarmuka pada koordinat lokal adalah sebagai berikut.
\begin{equation*}
u = 0
\end{equation*}
\begin{equation*}
E_{f}(0) = E_{c}(d_{c})
\end{equation*}
Kontinuitas $H_{y}$ menuntut turunan terbobot yang sama.
\begin{equation*}
\left.\frac{1}{\mu}\frac{dE_{f}}{du}\right|_{0}
=
\left.\frac{dE_{c}}{dz}\right|_{d_{c}}
\end{equation*}
Normalisasi pusat dipilih sebagai berikut.
\begin{equation*}
E_{c}(0) = 1
\end{equation*}
Penyelesaian sistem dua persamaan itu adalah sebagai berikut.
Catatan notasi: $\mathcal{A},\mathcal{B}$ (huruf kaligrafi) adalah koefisien
medan listrik, sedangkan $A,B$ pada bab permeabilitas adalah kekakuan magnetik.
\begin{subequations}\label{eq:5}
\begin{equation}
\mathcal{A} = \cosh(k_{c}d_{c})
\label{eq:5a}
\end{equation}
\begin{equation}
\mathcal{B} = \frac{\mu\,k_{c}\,\sinh(k_{c}d_{c})}{k_{f}}
\label{eq:5b}
\end{equation}
\end{subequations}
Inilah fungsi \texttt{zOne} pada kode: hitung $\chi$, lalu $\mu$,
lalu $k_{c}$ dan $k_{f}$ kompleks, lalu $\mathcal{A}$ dan $\mathcal{B}$.

\subsection{Dari medan permukaan ke impedansi strip}
Medan di permukaan luar adalah nilai solusi pada tebal penuh.
\begin{equation*}
E_{s} = E_{f}(t_{f})
\end{equation*}
Medan magnetik di permukaan luar adalah sebagai berikut.
\begin{subequations}\label{eq:6}
\begin{equation}
H_{s} = \frac{k_{f}}{j\omega\mu_{0}\mu}\,
\bigl[\mathcal{A}\sinh(k_{f}t_{f})+\mathcal{B}\cosh(k_{f}t_{f})\bigr]
\label{eq:6a}
\end{equation}
Hukum Ampere pada strip lebar, dengan $H_{s}$ di atas, memberi total arus sebagai berikut.
\begin{equation*}
I = 2\,w\,H_{s}
\end{equation*}
Tegangan sepanjang strip adalah sebagai berikut.
\begin{equation*}
V = E_{s}\,l
\end{equation*}
Impedansi adalah perbandingan keduanya sebagai berikut.
\begin{equation}
Z = \frac{l}{2w}\,\frac{E_{s}}{H_{s}}
\label{eq:6b}
\end{equation}
\end{subequations}
Faktor setengah muncul karena arus terbagi pada dua muka simetris.

\subsection{Batas DC dan uji kewajaran}
Pada frekuensi nol, bilangan gelombang nol dan impedansi harus menjadi
resistansi DC paralel dua lapis magnetik plus inti.
\begin{equation}
R_{\mathrm{dc}} =
\frac{l}{w\,(\sigma_{c}\,2d_{c}+2\sigma_{f}\,t_{f})}
\label{eq:7}
\end{equation}
Untuk satu bahan homogen, rumus umum menyusut ke bentuk klasik berikut.
\begin{equation*}
Z = R_{\mathrm{dc}}\,k\,t\coth(k\,t)
\end{equation*}
Pada frekuensi rendah, bentuk itu menuju resistansi DC sebagai berikut.
\begin{equation*}
Z \to R_{\mathrm{dc}}\quad\text{untuk}\quad f\to 0
\end{equation*}
Kedua batas ini dipakai sebagai uji numerik kewajaran simulator.

%=====================================================================
\section{Permeabilitas dari LLG dan Smit--Beljers}
%=====================================================================
\subsection{Persamaan gerak magnetisasi}
Vektor satuan magnetisasi $m$ berpresesi mengelilingi medan efektif
dan perlahan teredam menuju medan itu.
\begin{equation}
\frac{d\mathbf{m}}{dt}
= -\gamma\mu_{0}\,\mathbf{m}\times\mathbf{H}_{\mathrm{eff}}
+ \alpha\,\mathbf{m}\times\frac{d\mathbf{m}}{dt}
\label{eq:8}
\end{equation}
Suku pertama adalah presesi tanpa rugi.
Suku kedua adalah redaman Gilbert yang sebanding dengan $\alpha$.

\subsection{Ide linearisasi: dua kekakuan}
Gangguan kecil disekitar keseimbangan ditulis sebagai berikut.
\begin{equation*}
\mathbf{m} = \mathbf{m}_{0}+\delta\mathbf{m}\,e^{j\omega t}
\end{equation*}
Energi terhadap dua arah simpangan memberi dua kekakuan:
$A$ untuk simpangan dalam bidang,
$B$ untuk simpangan keluar bidang.
Suseptibilitas besar bila salah satu kekakuan kecil,
yaitu sistem lembek terhadap gangguan.

\subsection{Kasus longitudinal}
Sumbu mudah sejajar medan dan arus.
Komponen medan sepanjang magnetisasi adalah sebagai berikut
($s$ adalah cabang $+1$ atau $-1$).
\begin{equation*}
H_{p} = \max(s\,H,\,-0.95\,H_{k})
\end{equation*}
Batas bawah $-0.95H_{k}$ menjaga kekakuan tetap positif
tepat sebelum pembalikan.
Dua kekakuan longitudinal adalah sebagai berikut.
\begin{equation*}
A = H_{k}+H_{p}
\end{equation*}
\begin{equation*}
B = H_{k}+H_{p}+\mu_{0}M_{s}
\end{equation*}
Pada medan nol, kekakuan sama dengan anisotropi sebagai berikut.
\begin{equation*}
A = H_{k}
\end{equation*}
Nilai itu kecil sehingga suseptibilitas besar.
Inilah asal puncak tunggal runcing di nol.

\subsection{Kasus transversal}
Medan diambil nilai mutlaknya sebagai berikut.
\begin{equation*}
H = |H_{m}|
\end{equation*}
Bila medan di atas anisotropi, magnetisasi ikut medan.
\begin{equation*}
A = H-H_{k}\quad\text{untuk}\quad H\ge H_{k}
\end{equation*}
\begin{equation*}
B = H+\mu_{0}M_{s}\quad\text{untuk}\quad H\ge H_{k}
\end{equation*}
Bila medan di bawah anisotropi, magnetisasi miring dengan sudut yang
memenuhi kosinus berikut.
\begin{equation*}
\cos\phi = H/H_{k}
\end{equation*}
Kekakuan dalam bidang menjadi sebagai berikut.
\begin{equation*}
A = (H_{k}^{2}-H^{2})/H_{k}\quad\text{untuk}\quad H<H_{k}
\end{equation*}
Kekakuan luar bidang tetap sebagai berikut.
\begin{equation*}
B = H_{k}+\mu_{0}M_{s}\quad\text{untuk}\quad H<H_{k}
\end{equation*}
Hanya proyeksi sepanjang medan yang terukur sehingga muncul faktor reduksi.
\begin{equation*}
c_{2} = (H/H_{k})^{2}\quad\text{untuk}\quad H<H_{k}
\end{equation*}
\begin{equation*}
c_{2} = 1\quad\text{untuk}\quad H\ge H_{k}
\end{equation*}
Kekakuan dalam bidang nol pada resonansi sebagai berikut.
\begin{equation*}
A = 0
\end{equation*}
Kondisi resonansi terjadi pada medan berikut.
\begin{equation*}
H = H_{k}
\end{equation*}
Sehingga muncul dua puncak pada plus-minus anisotropi.

\subsection{Frekuensi resonansi, suseptibilitas, dan redaman}
Frekuensi resonansi feromagnetik adalah sebagai berikut.
\begin{equation}
f_{r} = \frac{\gamma}{2\pi}\sqrt{AB}
\label{eq:9}
\end{equation}
Suseptibilitas kompleks hasil linearisasi adalah sebagai berikut.
\begin{subequations}\label{eq:10}
\begin{equation}
\chi = c_{2}\,\mu_{0}M_{s}\,
\frac{B+j\Gamma}{(A+j\Gamma)(B+j\Gamma)-(\omega/\gamma)^{2}}
\label{eq:10a}
\end{equation}
Permeabilitas relatif adalah sebagai berikut.
\begin{equation}
\mu = 1+\chi
\label{eq:10b}
\end{equation}
\end{subequations}
Lebar garis total terdiri dari redaman Gilbert plus sebaran inhomogen.
\begin{equation}
\Gamma = \alpha\,\omega/\gamma+\Delta H_{0}/2
\label{eq:11}
\end{equation}
Jauh di bawah resonansi, penyebut didominasi $AB$ sehingga berlaku limit
statis sebagai berikut.
\begin{equation*}
\chi \to \mu_{0}M_{s}/A
\end{equation*}
Khusus longitudinal, limitnya adalah sebagai berikut.
\begin{equation*}
\chi_{0} = \mu_{0}M_{s}/(H+H_{k})
\end{equation*}

\subsection{Koreksi tepi film berhingga}
Penampang lintang lapisan magnetik berupa lempeng tipis setebal $t_{f}$
dan selebar $w$. Pendekatan penampang elips memberi faktor demagnetisasi
sepanjang lebar sebagai berikut.
\begin{equation*}
N_{w} = \frac{t_{f}}{t_{f}+w}
\end{equation*}
Anisotropi bentuk bersifat \emph{aditif} pada medan anisotropi.
Untuk sumbu mudah longitudinal, demagnetisasi sepanjang lebar menaikkan
kekakuan; untuk transversal, ia menurunkannya.
\begin{equation*}
H_{k,\mathrm{eff}} = H_{k} \pm \mathrm{edgeC}\,N_{w}\,\mu_{0}M_{s}
\end{equation*}
Tanda plus untuk longitudinal dan minus untuk transversal.
Hanya tebal lapisan magnetik yang ikut; tembaga tidak.
Nilai nol berarti mati (perilaku tanpa koreksi) dan nilai satu berarti
koreksi penuh. Satuan: $t_{f}$ dalam mikrometer, $w$ dalam milimeter
dikali seribu. Bentuk ini mengganti bentuk perkalian lama,
yang tidak dapat dibenarkan secara fisis.

\subsection{Dinding domain empiris frekuensi rendah}
Gerak dinding domain menambah suseptibilitas real pada frekuensi rendah.
\begin{equation*}
\chi_{\mathrm{dw}}(f) =
\frac{\mathrm{dwAmp}}{100}\,\frac{\mu_{0}M_{s}}{A_{\ast}}\,
\frac{1}{1+j\,f/f_{\mathrm{dw}}},\qquad f_{\mathrm{dw}}=5\ \mathrm{MHz}
\end{equation*}
Ini bentuk Debye. Bagian real dan imajiner ditambahkan bersama:
\begin{equation*}
\Delta\chi' = \chi_{0}/(1+u^{2}),\qquad
\Delta\chi'' = -\chi_{0}\,u/(1+u^{2}),\qquad u=f/f_{\mathrm{dw}}
\end{equation*}
Bentuk ini memenuhi relasi Kramers--Kronig. Versi sebelumnya hanya menambah
bagian real, yang melanggar kausalitas. Frekuensi $5$ MHz adalah pilihan empiris.
Penyebut diamankan agar tidak divergen sebagai berikut.
\begin{equation*}
A_{\ast} = \max(A,\,0.05\,H_{k},\,10^{-9})
\end{equation*}
Tanpa pengaman ini, titik nol berikut pada resonansi transversal
akan menghasilkan nilai tak-hingga.
\begin{equation*}
A = 0
\end{equation*}
Nilai nol berarti mati.

\subsection{Histeresis koersif sebagai campuran dua keadaan}
Magnetisasi longitudinal berbalik di sekitar plus-minus $H_{c}$.
Alih-alih loncatan tajam, dipakai fungsi alih mulus berupa
fungsi distribusi kumulatif normal $\Phi$.
Peluang keadaan plus pada sapuan naik adalah sebagai berikut.
\begin{equation*}
p_{\mathrm{up}} = \Phi\!\left(\frac{H-H_{c}}{\Delta_{c}}\right)
\end{equation*}
Peluang keadaan plus pada sapuan turun adalah sebagai berikut.
\begin{equation*}
p_{\mathrm{down}} = \Phi\!\left(\frac{H+H_{c}}{\Delta_{c}}\right)
\end{equation*}
Batas koersif dijaga di bawah anisotropi sebagai berikut.
\begin{equation*}
H_{c} = \min(H_{c},\,0.95\,H_{k})
\end{equation*}
Lebar alih $\Delta_{c}$ (bukan konduktivitas) terdiri dari bawaan plus tambahan pengguna.
\begin{equation*}
\Delta_{c} = \max(0.3\,H_{c},\,0.03)+\mathrm{sigHc}
\end{equation*}
Impedansi campuran adalah rata-rata berbobot dua keadaan.
\begin{equation*}
Z = p\,Z_{+}+(1-p)\,Z_{-}
\end{equation*}
Bila $p$ jenuh di atas $0.9995$ atau di bawah $0.0005$,
kode langsung memakai satu cabang agar hemat.

\subsection{Sebaran anisotropi film}
Film nyata tidak seragam.
Rata-rata atas sebaran normal $H_{k}$ memakai kuadratur tujuh titik.
\begin{equation*}
\bar Z(H) = \sum_{i=1}^{7} w_{i}\,Z(H;\,H_{k}+\sigma_{Hk}\,x_{i})
\end{equation*}
Titik yang keluar batas dijaga agar tetap positif sebagai berikut.
\begin{equation*}
H_{k,i} = \max(H_{k}+\sigma_{Hk}\,x_{i},\,0.05)
\end{equation*}
Nilai nol berarti seragam dan sama dengan perilaku lama.
Bobot dan simpul $x_{i},w_{i}$ ada pada lampiran.

%=====================================================================
\section{Respons nanopartikel}
%=====================================================================
Volume bola dari diameter $d$ adalah sebagai berikut.
\begin{equation*}
V = \pi d^{3}/6
\end{equation*}
Magnetisasi dalam A/m dari masukan kilo-A/m adalah sebagai berikut.
\begin{equation*}
M_{s}[\mathrm{A/m}] = M_{sp}[\mathrm{kA/m}]\times 10^{3}
\end{equation*}

\subsection{Superparamagnetik: Langevin dan relaksasi Neel}
Momen satu partikel adalah sebagai berikut.
\begin{equation*}
m_{p} = M_{s}\,V
\end{equation*}
Parameter perbandingan energi Zeeman terhadap termal adalah sebagai berikut.
\begin{subequations}\label{eq:12}
\begin{equation}
\xi = \frac{m_{p}\,B}{k_{B}T}
\label{eq:12a}
\end{equation}
Fungsi Langevin didefinisikan sebagai berikut.
\begin{equation}
L(\xi) = \coth\xi-1/\xi
\label{eq:12b}
\end{equation}
Magnetisasi rata-rata mengikuti fungsi itu sebagai berikut.
\begin{equation}
M = M_{s}\,L(\xi)
\label{eq:12c}
\end{equation}
\end{subequations}
Untuk $\xi$ kecil dipakai deret agar stabil sebagai berikut.
\begin{equation*}
L(\xi) \approx \xi/3-\xi^{3}/45
\end{equation*}
Suseptibilitas transversal adalah magnetisasi dibagi medan.
\begin{subequations}\label{eq:13}
\begin{equation}
\chi_{\perp} = \frac{\mu_{0}M_{s}\,L(\xi)}{B}
\label{eq:13a}
\end{equation}
Limit medan nol memberi bentuk Curie sebagai berikut.
\begin{equation}
\chi(0) = \frac{\mu_{0}M_{s}^{2}\,V}{3k_{B}T}
\label{eq:13b}
\end{equation}
\end{subequations}
Waktu Neel untuk lompatan termal adalah sebagai berikut.
\begin{subequations}\label{eq:14}
\begin{equation}
\tau_{N} = \tau_{0}\,\exp(KV/k_{B}T)
\label{eq:14a}
\end{equation}
Eksponen dibatasi sampai 60 agar tidak meluap.
Respons Debye pada frekuensi ukur adalah sebagai berikut.
\begin{equation}
\chi(\omega) = \frac{\chi_{\perp}}{1+j\omega\tau_{N}}
\label{eq:14b}
\end{equation}
\end{subequations}
Bentuk real dan imajiner yang dipakai kode adalah sebagai berikut
($u$ adalah frekuensi ternormalisasi).
\begin{equation*}
u = \omega\,\tau_{N}
\end{equation*}
\begin{equation*}
\mathrm{Re}\,\chi = \chi_{0}/(1+u^{2})
\end{equation*}
\begin{equation*}
\mathrm{Im}\,\chi = -\chi_{0}\,u/(1+u^{2})
\end{equation*}

\subsection{Terblokir: Stoner--Wohlfarth acak plus resonansi}
Suseptibilitas awal ansambel acak adalah sebagai berikut.
\begin{subequations}\label{eq:15}
\begin{equation}
\chi_{i} = \frac{\mu_{0}M_{s}^{2}}{3K}
\label{eq:15a}
\end{equation}
Medan anisotropi partikel adalah sebagai berikut.
\begin{equation}
B_{K} = \frac{2K}{M_{s}}
\label{eq:15b}
\end{equation}
Medan koersif rata-rata ansambel acak adalah sebagai berikut.
\begin{equation}
B_{c} \approx 0.479\,B_{K}
\label{eq:15c}
\end{equation}
\end{subequations}
Magnetisasi statis didekati tangen hiperbolik yang digeser histeresis.
\begin{equation*}
s_{h} = B\mp B_{c}
\end{equation*}
Tanda atas untuk sapuan naik, tanda bawah untuk sapuan turun.
\begin{equation*}
a = \chi_{i}\,s_{h}/(\mu_{0}M_{s})
\end{equation*}
\begin{subequations}\label{eq:16}
\begin{equation}
M = M_{s}\,\tanh(a)
\label{eq:16a}
\end{equation}
Medan RF sensor tegak lurus terhadap medan bias, sehingga yang relevan
untuk kopling adalah suseptibilitas \emph{transversal}, bukan turunan
magnetisasi. Dipakai bentuk interpolasi berikut, yang bernilai
$\chi_{i}$ di medan nol dan menuju limit sekan superparamagnetik
$\mu_{0}M_{s}/B$ pada medan tinggi.
\begin{equation*}
\chi_{\perp,b} = \frac{\chi_{i}}{1+|B|/(1.5\,B_{K})}
\end{equation*}
Histeresis hanya masuk lewat magnetisasi $M$. Bentuk ini heuristik.
Versi sebelumnya memakai turunan $\chi_{i}/\cosh^{2}(a)$,
yang tidak konsisten dengan model superparamagnetik.
Frekuensi resonansi partikel adalah sebagai berikut.
\begin{equation}
\omega_{r} = \gamma\,(B_{K}+B)
\label{eq:16b}
\end{equation}
Respons dinamik berbentuk osilator Lorentz sebagai berikut.
\begin{equation}
\chi(\omega) =
\frac{\chi_{\perp,b}}{1-(\omega/\omega_{r})^{2}+2j\alpha_{p}\,\omega/\omega_{r}}
\label{eq:16c}
\end{equation}
\end{subequations}
Simbol $\alpha_{p}$ adalah redaman partikel (masukan panel, tipikal $0.05$).
Faktor amplitudo momen untuk sebaran dipolar adalah sebagai berikut
($\varpi$ adalah magnetisasi ternormalisasi).
\begin{equation*}
\varpi = M/M_{s}
\end{equation*}
\begin{equation*}
s = \sqrt{\varpi^{2}+(1-\varpi^{2})/3}
\end{equation*}
Bentuk itu adalah interpolasi heuristik untuk akar rata-rata kuadrat
proyeksi momen. Saturasi penuh memberi kuadrat satu sebagai berikut.
\begin{equation*}
s^{2} = 1
\end{equation*}
Acak penuh memberi kuadrat sepertiga sebagai berikut.
\begin{equation*}
s^{2} = 1/3
\end{equation*}
Sehingga faktornya sendiri adalah satu per akar tiga sebagai berikut.
\begin{equation*}
s = 1/\sqrt{3}\approx 0.577
\end{equation*}
Untuk partikel superparamagnetik, bentuk eksaknya adalah
\begin{equation*}
\langle\cos^{2}\theta\rangle = 1-2L(\xi)/\xi
\end{equation*}
dan bentuk itu tidak dipakai pada kode.

\subsection{Paramagnetik, diamagnetik, dan rugi eddy logam}
Untuk tipe para dan dia, suseptibilitas volume dianggap konstan.
\begin{equation*}
\chi = \chi_{V}
\end{equation*}
\begin{equation*}
M_{p} = \chi_{V}\,B/\mu_{0}
\end{equation*}
Kode menyimpan $\mu_{0}M_{p}$ (dalam tesla), yaitu $\chi_{V}B$,
sehingga $B_{s}$ pada \eqref{eq:20} konsisten dimensinya.
\begin{equation*}
s = 0
\end{equation*}
Bola logam berarus eddy menambah bagian imajiner negatif.
Pendekatan frekuensi rendah diturunkan dari disipasi daya. Pada medan
uniform berosilasi, medan induksi azimut di dalam bola berjari-jari $a$
adalah sebagai berikut.
\begin{equation*}
E_{\varphi} = -j\omega B_{0}\,r\sin\theta/2
\end{equation*}
Daya rata-rata volume bola adalah sebagai berikut.
\begin{equation*}
P = \frac{\sigma}{2}\int |E_{\varphi}|^{2}\,dV
= \frac{\pi}{15}\,\sigma\,\omega^{2}B_{0}^{2}\,a^{5}
\end{equation*}
Daya yang sama harus sama dengan serapan oleh suseptibilitas imajiner.
Volume bola adalah sebagai berikut.
\begin{equation*}
V = 4\pi a^{3}/3
\end{equation*}
\begin{equation*}
P = \frac{\omega}{2}\,\frac{\chi''B_{0}^{2}}{\mu_{0}}\,V
\end{equation*}
Menyamakan keduanya memberi bentuk yang dipakai.
\begin{equation*}
\chi_{\mathrm{eddy}} = -j\,\mu_{0}\,\sigma\,\omega\,(d/2)^{2}/10
\end{equation*}
Bentuk ekuivalen dengan diameter penuh adalah sebagai berikut.
\begin{equation*}
\chi_{\mathrm{eddy}} = -j\,\mu_{0}\,\sigma\,\omega\,d^{2}/40
\end{equation*}
Suku ini hanya aktif bila konduktivitas partikel positif.
Nilai nol berarti isolator dan perilaku lama dipertahankan.

\subsection{Sebaran diameter lognormal yang menjaga rata-rata}
Diameter mengikuti lognormal dengan simpangan log $\sigma_{D}$.
Diameter $d$ pada panel adalah rata-rata berbobot jumlah.
Titik kuadratur yang menjaga mean adalah sebagai berikut.
\begin{equation*}
d_{i} = d\,\exp(\sigma_{D}\,x_{i}-\sigma_{D}^{2}/2)
\end{equation*}
Fraksi volume $\varphi$ dan massa adalah besaran volume, sehingga rata-rata
besaran $q$ (magnetisasi, faktor $s$, dan $\chi$) memakai bobot volume.
\begin{equation*}
\bar q = \frac{\sum_{i} w_{i}\,d_{i}^{3}\,q(d_{i})}{\sum_{i} w_{i}\,d_{i}^{3}}
\end{equation*}
Versi sebelumnya memakai bobot jumlah partikel yang tidak konsisten dengan
$\varphi$.
Nilai nol berarti monodispers dan sama dengan satu pemanggilan langsung.

%=====================================================================
\section{Kopling lapisan partikel ke sensor}
%=====================================================================
\subsection{Tebal lapisan dan kedalaman interaksi}
Massa $m$, massa jenis $\rho$, fraksi volume $\varphi$,
luas sebaran $A_{c}$ memberi tebal sebagai berikut.
\begin{subequations}\label{eq:17}
\begin{equation}
t_{p} = \frac{m}{\rho\,\varphi\,A_{c}}
\label{eq:17a}
\end{equation}
Skala interaksi diambil sebanding dengan lebar strip. Ini \textbf{heuristik},
bukan turunan: medan sebar strip sebenarnya meluruh secara aljabar dan
tidak bergantung frekuensi, bukan secara eksponensial.
\begin{equation}
\lambda = \frac{w}{\pi}
\label{eq:17b}
\end{equation}
Hanya lapisan setebal orde $\lambda$ yang berinteraksi,
sehingga dipakai tebal efektif jenuh sebagai berikut.
\begin{equation}
t_{\mathrm{eff}} = \lambda\,\bigl(1-e^{-t_{p}/\lambda}\bigr)
\label{eq:17c}
\end{equation}
\end{subequations}
Bila lapisan tipis, tebal efektif sama dengan tebal nyata.
Bila lapisan tebal, tebal efektif jenuh di $\lambda$.
Inilah alasan kurva kalibrasi massa melengkung dan jenuh.

\subsection{Faktor demagnetisasi spheroid pipih}
Jari-jari ekuivalen sebaran adalah sebagai berikut.
\begin{equation*}
a = \sqrt{A_{c}/\pi}
\end{equation*}
Rasio sumbu pipih adalah sebagai berikut.
\begin{equation*}
e = (t_{p}/2)/a
\end{equation*}
Faktor demagnetisasi tegak lurus adalah sebagai berikut.
\begin{equation*}
N_{z} = \frac{1-e\,\arccos e/\sqrt{1-e^{2}}}{1-e^{2}}
\end{equation*}
Kasus batas ditangani sebagai berikut.
\begin{equation*}
N_{z} = 1/3\quad\text{untuk}\quad e\ge 1
\end{equation*}
\begin{equation*}
N_{z} = 1\quad\text{untuk}\quad e\le 0
\end{equation*}
Faktor dalam bidang mengikuti jejak tensor satuan.
\begin{equation*}
N_{\mathrm{in}} = (1-N_{z})/2
\end{equation*}
Jadi film sangat tipis memberi faktor dalam bidang menuju nol,
dan bola memberi sepertiga, sebagai berikut.
\begin{equation*}
N_{\mathrm{in}}\to 0\quad\text{film tipis},\qquad
N_{\mathrm{in}} = 1/3\quad\text{bola}
\end{equation*}
Fungsi \texttt{Nin} pada kode langsung
mengembalikan $N_{\mathrm{in}}$ dengan nilai batas tersebut.
Rumus arccos hanya sahih untuk $e<1$; untuk $e\ge 1$ dipakai nilai bola
sebagai pendekatan.

\subsection{Mekanisme 1: kopling dinamis Maxwell--Garnett}
Campuran partikel dalam matriks dianggap medium efektif.
Fraksi yang dipakai adalah fraksi volume $\varphi$.
\begin{subequations}\label{eq:18}
\begin{equation}
\frac{\mu_{L}-1}{\mu_{L}+2}
=
\varphi\,\frac{\mu_{p}-1}{\mu_{p}+2}
\label{eq:18a}
\end{equation}
Permeabilitas partikel adalah sebagai berikut.
\begin{equation}
\mu_{p} = 1+\chi
\label{eq:18b}
\end{equation}
\end{subequations}
Bentuk pecahan di atas diselesaikan untuk $\mu_{L}$ pada kode
dengan aritmetika kompleks eksplisit.
Lapisan itu menambah fluks pada satu muka strip sehingga muncul
tambahan impedansi seri sebagai berikut.
\begin{equation}
\Delta Z_{p} = j\,\omega\,\mu_{0}\,(\mu_{L}-1)\,t_{\mathrm{eff}}\,l/(2w)
\label{eq:19}
\end{equation}
Bagian real $\Delta Z_{p}$ berasal dari bagian imajiner $\mu_{L}$,
bagian imajiner $\Delta Z_{p}$ berasal dari bagian real $\mu_{L}$.

\subsection{Mekanisme 2: medan sebar rata-rata}
Lapisan termagnetisasi menimbulkan medan demagnetisasi melawan medan luar.
Komponen tangensial kontinu sehingga sensor merasakan medan efektif
lebih kecil sebagai berikut.
\begin{equation}
B_{s} = -N_{\mathrm{in}}\,\varphi\,\mu_{0}\,M_{p}(H)
\label{eq:20}
\end{equation}
Medan efektif pada sensor adalah jumlah medan luar dan sebaran.
\begin{equation*}
H_{\mathrm{efektif}} = H_{m}+B_{s}\times 10^{3}
\end{equation*}
Faktor seribu adalah konversi tesla ke mili-tesla.

\subsection{Mekanisme 3: sebaran dipolar dan rata-rata medan lokal}
Medan dipol tetangga berfluktuasi dari titik ke titik.
Lebar sebaran diperkirakan dari orde medan dipol sebagai berikut.
\begin{subequations}\label{eq:21}
\begin{equation}
\sigma_{s} = C_{d}\,\varphi\,\mu_{0}\,M_{s}\,s
\label{eq:21a}
\end{equation}
Faktor $C_{d}$ adalah koreksi perataan oleh kopling tukar dan geometri,
suku paling tak-pasti pada model.
Sensor merespons rata-rata atas sebaran itu sebagai berikut.
\begin{equation}
\bar Z = \sum_{i} w_{i}\,Z(H+B_{s}+\sigma_{s}\,x_{i})
\label{eq:21b}
\end{equation}
\end{subequations}
Rata-rata dihitung dengan kuadratur tujuh titik yang sama.

\subsection{Impedansi total, cakupan, dua sisi, dan kabel}
Jumlah semua kontribusi adalah sebagai berikut.
\begin{equation}
Z_{\mathrm{tot}} = \bar Z+\Delta Z_{p}+R_{c}+j\,\omega\,L_{p}
\label{eq:22}
\end{equation}
Cakupan pulau $\eta$ mengecilkan ketiga efek partikel.
\begin{equation*}
B_{s} \to \eta\,B_{s}
\end{equation*}
\begin{equation*}
\sigma_{s} \to \eta\,\sigma_{s}
\end{equation*}
\begin{equation*}
\Delta Z_{p} \to \eta\,\Delta Z_{p}
\end{equation*}
Nilai satu berarti film kontinu dan sama dengan perilaku lama.
Nilai nol berarti tidak ada interaksi.
Lapisan dua sisi menggandakan kopling dinamis sebagai berikut.
\begin{equation*}
\Delta Z_{p} \to 2\,\Delta Z_{p}
\end{equation*}
Resistansi kabel termasuk efek kulit sebagai berikut.
\begin{equation*}
R_{c}(f) = R_{c0}+R_{c,\mathrm{skin}}\,\sqrt{f_{\mathrm{MHz}}}
\end{equation*}
Arah perubahan puncak dapat dipahami sebagai berikut.
Mekanisme dinamis cenderung menaikkan puncak bila $\chi$ besar di nol
(khas superparamagnetik).
Mekanisme sebaran cenderung menurunkan dan melebarkan puncak
(khas terblokir bermomen beku).
Simulator menampilkan hasil bersih apa adanya.

%=====================================================================
\section{Sapuan, derau, pita model, dan batas deteksi}
%=====================================================================
\subsection{Titik sapuan dengan zag halus di tengah}
Himpunan titik positif dibangun dari langkah halus lalu kasar,
selalu mencakup nol dan ujung sebagai berikut.
\begin{equation*}
\mathcal{P} = \{0\}\cup\{d_{1},2d_{1},\dots\}\cup\{\dots\}\cup\{H_{\max}\}
\end{equation*}
Cabang negatif adalah cerminannya, lalu digabung terurut.
Sapuan naik memakai urutan itu, sapuan turun memakai kebalikannya.

\subsection{Model tanpa derau pada tiap titik}
Acuan tiap massa adalah modulus pada ujung naik sebagai berikut.
\begin{equation*}
Z_{\mathrm{ref}} = |Z(H_{\max}^{\mathrm{akt}},\mathrm{naik})|
\end{equation*}
Kurva model tanpa derau pada langkah halus adalah sebagai berikut.
\begin{equation*}
y(H) = 100\,(|Z(H)|-Z_{\mathrm{ref}})/Z_{\mathrm{ref}}
\end{equation*}
Puncak, posisi, dan lebar setengah-tinggi dibaca dari kurva itu
dengan interpolasi linear antar titik.

\subsection{Data terukur dengan derau ganda}
Acuan terukur ikut berderau sebagai berikut.
\begin{equation*}
Z_{\mathrm{ref,meas}} = Z_{\mathrm{ref}}\,(1+\sigma_{Z}\,n_{0}/100)
\end{equation*}
Tiap titik memakai medan aktual fisis lalu dibaca dengan derau probe.
Modulus ikut berderau sebagai berikut.
\begin{equation*}
Z_{\mathrm{meas}} = |Z(H_{\mathrm{act}})|\,(1+\sigma_{Z}\,n_{1}/100)
\end{equation*}
MIR terukur memakai acuan terukur yang sama sebagai berikut.
\begin{equation*}
\mathrm{MIR}_{\mathrm{meas}} = 100\,(Z_{\mathrm{meas}}-Z_{\mathrm{ref,meas}})/Z_{\mathrm{ref,meas}}
\end{equation*}
Derau $n$ adalah normal baku dari generator deterministik agar
dapat diulang dengan benih yang sama.

\subsection{Sensitivitas awal dan batas deteksi}
Sensitivitas adalah turunan puncak terhadap massa di nol,
didekati beda hingga sebagai berikut. Langkah dipilih adaptif
agar tetap pada rezim linear, yaitu pada lima persen lambda sebagai berikut.
\begin{equation*}
t_{p} = 0.05\,\lambda
\end{equation*}
Nilai tetap $0.02$ g hanya dipakai bila tidak ada partikel.
\begin{equation*}
dm = 0.05\,\lambda\,\rho\,\varphi\,A_{c}
\end{equation*}
\begin{equation*}
S = \bigl[P(dm)-P(0)\bigr]/dm
\end{equation*}
Derau MIR pada puncak mengikuti propagasi dua pembacaan berderau.
\begin{subequations}\label{eq:23}
\begin{equation}
\sigma_{\mathrm{MIR}} \approx (1+\mathrm{MIR}/100)\,\sqrt{2}\,\sigma_{Z}
\label{eq:23a}
\end{equation}
Batas deteksi tiga-sigma adalah sebagai berikut.
\begin{equation}
m_{\mathrm{LOD}} = 3\,\sigma_{\mathrm{MIR}}/|S|
\label{eq:23b}
\end{equation}
\end{subequations}
Bila sensitivitas nol, batas menjadi tak-hingga.

\subsection{Satu definisi sigma untuk pita, residu, dan fit}
Derau medan dipropagasi lewat kemiringan kurva.
Gabungan derau aktual dan derau baca adalah sebagai berikut.
\begin{equation*}
\sigma_{H,\mathrm{tot}} = \sqrt{\sigma_{H}^{2}+\sigma_{Hr}^{2}}
\end{equation*}
Suku derau impedansi dalam poin persen adalah sebagai berikut.
\begin{equation*}
\sigma_{Z}^{\prime} = \sigma_{Z}\,\sqrt{2}\,(1+m/100)
\end{equation*}
Gabungan total adalah sebagai berikut.
\begin{equation*}
\sigma = \sqrt{(|d\mathrm{MIR}/dH|\,\sigma_{H,\mathrm{tot}})^{2}
+{\sigma_{Z}^{\prime}}^{2}}
\end{equation*}
Bawahnya dibatasi sangat kecil agar tidak nol persis.
Kemiringan memakai stencil lima titik agar tidak bergerigi di puncak tajam.
Pita model adalah kurva plus-minus satu sigma sebagai berikut.
\begin{equation*}
y_{\mathrm{atas}} = y+\sigma
\end{equation*}
\begin{equation*}
y_{\mathrm{bawah}} = y-\sigma
\end{equation*}
Residu eksperimen adalah selisih terhadap model sebagai berikut.
\begin{equation*}
r = \mathrm{MIR}_{\mathrm{exp}}-\mathrm{MIR}_{\mathrm{mod}}
\end{equation*}
Kecocokan per derajat bebas mengikuti bentuk khi-kuadrat sebagai berikut.
\begin{equation*}
\chi^{2}/\mathrm{DoF} = \frac{1}{\mathrm{DoF}}\sum (r/\sigma)^{2}
\end{equation*}
Nilai sekitar satu berarti konsisten, jauh di atas satu berarti salah model
sistematis, bukan sekadar kurang derau.
Catatan: derau acuan $n_{0}$ dipakai bersama oleh semua titik satu sampel
(korelasi common-mode), sedangkan $\sigma$ di atas menganggap titik
independen, sehingga $\chi^{2}/\mathrm{DoF}$ sedikit diremehkan.
Model untuk data campuran arah memakai rata-rata naik-turun,
yang menghapus histeresis.
\begin{equation*}
\mathrm{MIR}_{\mathrm{mod}}(H) =
100\,\bigl(0.5\,(|Z_{u}|+|Z_{d}|)-Z_{\mathrm{ref}}\bigr)/Z_{\mathrm{ref}}
\end{equation*}
Optimasi tanpa turunan memakai simpleks Nelder--Mead:
pantulkan, kembangkan, kerutkan, dan susutkan sampai sebaran nilai kecil.
Protokol yang disarankan adalah kunci sensor dari sampel nol dulu,
baru cocokkan kopling global, dan uji buta pada massa yang tidak ikut fit.

\appendix
\section{Tabel kuadratur tujuh titik}
Nilai berikut adalah simpul normal baku dan bobotnya.
Jumlah bobot adalah satu.
\begin{center}
\begin{tabular}{cc}
\toprule
Simpul & Bobot\\
\midrule
0 & 0.457143\\
plus-minus 1.154405 & 0.240123\\
plus-minus 2.366759 & 0.0307571\\
plus-minus 3.750440 & 0.000548269\\
\bottomrule
\end{tabular}
\end{center}
Asal nilai adalah simpul Gauss--Hermite fisikawan yang \emph{dikali}
akar dua menjadi simpul probabilis sebagai berikut.
\begin{equation*}
x_{i} = \sqrt{2}\,y_{i}
\end{equation*}
Bobot fisikawan dibagi akar pi sebagai berikut.
\begin{equation*}
w_{i} = w_{i}^{\mathrm{fis}}/\sqrt{\pi}
\end{equation*}
Nilai dipakai untuk sebaran dipolar, sebaran anisotropi, dan diameter.
Rata-rata umum berbentuk jumlah berbobot sebagai berikut.
\begin{equation*}
\bar q = \sum_{i=1}^{7} w_{i}\,q(x_{i})
\end{equation*}

\section{Asumsi, batas, dan daftar cek konfirmasi}
Strip dianggap planar tak-hingga arah lebar sebagai berikut.
\begin{equation*}
w \gg t_{f},\,t_{\mathrm{cu}}
\end{equation*}
Implikasinya, efek tepi dan arus arah lebar diabaikan kecuali koreksi
orde pertama melalui parameter tepi.
Magnetisasi dianggap seragam berotasi koheren.
Gerak dinding domain dan resonansi ketebalan tukar tidak dimodelkan
penuh, sehingga sahih terutama di atas beberapa MHz pada film sederhana.
Arus perpindahan diabaikan.
Sensor dianggap terbias statis sejajar sumbu panjang.
Lapisan partikel dianggap medium efektif isolator.
Rugi eddy serbuk logam hanya lewat suseptibilitas kompleks,
sehingga hasil serbuk logam adalah perkiraan.
Kopling dinamis adalah pendekatan fluks satu muka, bukan solusi 3D penuh.
Semua parameter bahan adalah tipikal literatur, bukan kalibrasi.
Keterbatasan struktural yang diketahui dirinci pada Lampiran~\ref{app:empiris}.

Daftar cek dan statusnya:
batas DC dan batas homogen (otomatis, lulus),
simetri dan pembalikan waktu (otomatis, lulus),
koreksi tepi, dinding domain, $N_{\mathrm{in}}$, eddy, dan $\chi$ partikel
(otomatis, lulus), saturasi acuan di ujung (peringatan di antarmuka),
round-trip CSV dan uji buta pada massa yang tidak ikut fit
(belum otomatis, perlu data).


\section{Parameter empiris dan keterbatasan yang diketahui}\label{app:empiris}
Tabel berikut memisahkan suku yang diturunkan dari suku yang heuristik
atau berparameter bebas.
\begin{center}
\small
\begin{tabular}{p{0.30\textwidth}p{0.62\textwidth}}
\toprule
Suku & Status\\
\midrule
Skala panjang interaksi & Heuristik; lihat definisi lambda pada bab kopling.\\
$C_{d}$ & Parameter bebas; paling tak pasti pada model.\\
edgeC & Parameter bebas; bentuk aditif (lihat bab koreksi tepi).\\
dwAmp dan frekuensi relaksasi 5 MHz & Empiris; Debye kausal.\\
sigHc dan lebar bawaan tiga persepuluh koersif & Empiris; lebar peralihan histeresis.\\
Medan koersif ansatz dan tanh & Heuristik; remanen sekitar tiga persepuluh,
bukan setengah seperti Stoner--Wohlfarth murni.\\
Faktor amplitudo momen & Interpolasi heuristik
(lihat sub-bab faktor amplitudo).\\
$\chi_{\perp,b}$ & Interpolasi heuristik antara $\chi_{i}$ dan $\mu_{0}M_{s}/B$.\\
$\tau_{N}$ pada seluruh respons & Waktu lompat penghalang dipakai untuk respons
transversal; pendekatan.\\
Sebaran $H_{k}$, diameter, dipolar & Kuadratur Gaussian; sebaran dipolar
sebenarnya berekor berat (mendekati Lorentzian) sehingga ekornya diremehkan.\\
$B_{s}$ skalar seragam & Medan sebar lapisan berhingga tidak seragam dan dapat
berbalik tanda di dekat tepi.\\
Pemilihan SPM/terblokir & Manual; peralihan sekitar $KV/k_{B}T\approx 25$.\\
\bottomrule
\end{tabular}
\end{center}
Prediksi kuantitatif batas deteksi bergantung kuat pada suku-suku ini,
sehingga harus dikalibrasi terhadap data eksperimen.

\section{Hasil uji mandiri}\label{app:uji}
Uji berikut dijalankan pada kode simulator versi revisi
(sensor NiFe; anisotropi 3 mT; film 1 $\mu$m; tembaga 2 $\mu$m;
panjang 5 mm; lebar 0.5 mm). Semua lulus.
\begin{center}
\small
\begin{tabular}{p{0.46\textwidth}p{0.42\textwidth}}
\toprule
Uji & Hasil\\
\midrule
Batas DC, $Z(100\ \mathrm{Hz})\to R_{\mathrm{dc}}$ &
$0.0938087$ vs $0.0938086$ $\Omega$\\
Batas homogen, $R_{\mathrm{dc}}kt\coth kt$ (50 MHz) & galat relatif $9\times10^{-18}$\\
Simetri medan plus-minus & galat $0$\\
Pembalikan waktu dengan koersif & galat $0$\\
Koreksi tepi aditif & suseptibilitas statis $200.160$ (sama dengan teori)\\
Dinding domain Debye & beda real $16.667$, beda imajiner $-16.667$\\
Faktor dalam bidang film tipis & $1.6\times10^{-10}$\\
Faktor dalam bidang menuju bola & $0.3332$ (menuju sepertiga)\\
Rugi eddy & $-9.870\times10^{-5}$ (sama dengan teori seperempat-puluh)\\
Suseptibilitas transversal di medan nol & $0.25656$ (sama dengan awal)\\
Suseptibilitas transversal di medan seratus kali anisotropi &
rasio $0.9852$ (mendekati satu)\\
Sebaran diameter $\sigma_{D}\to0$ & galat $10^{-11}$\\
\bottomrule
\end{tabular}
\end{center}
Uji ini memeriksa konsistensi internal dan batas analitik.
Uji itu tidak menggantikan validasi terhadap data eksperimen.

\end{document}
